\section{Prototyping}
This section will explain how important prototyping and sketching was in the project and how it helped in the preanalysis. \\

Prototypes and sketches were used to quickly communicate and refine design ideas within the \localgroup{}, and for gathering feedback from customers. 
Many sketches and prototypes were produced over the course of the project, to keep communication about design effective. 
This also worked well with the iterative work form of the project, and new prototypes and sketches were continuously built on top of existing ones. \newline
Late prototypes were used to demonstrate the design of the product to the customers. 
These prototypes were implemented on paper, and showcased the different screens of the product while demonstrating the control flow of the system in specific use cases. 
This was important in evaluating the design, and gave reassurance about the quality of the design. \newline
An example of a prototype can be seen in \autoref{fig:first-prototype}

\begin{figure}[h!]
	\centering
	\includegraphics[scale=0.5]{gfx/first-prototype.pdf}
	\caption{Activity overview page from one of the first prototypes.}
	\label{fig:first-prototype}
\end{figure}